% !TeX root = ../main.tex

\chapter{Convex Optimization}
\label{chap:basics_co}

In this chapter, we will introduce the basics of convex optimization problems.
It will be useful for solving POVMs for quantum state discrimination.

\section{Basic terminology}
What is a convex programming problem?

To solve the formulated semidefinite programming problems, we take advantage of
the state-of-the-art framework CVXPY and solver SCS with acceleration
techniques applied.

We have also considered using other solvers like Gurobi and IBM CPlex, but due to limited
time, we haven't implemented all the programming problems for them.

\section{CVXPY}
CVXPY is an open source Python module for modeling convex optimization problems \cite{RN473}.



\subsection{Disciplined convex programming (DCP)}
Disciplined convex programming (DCP) is a system for constructing mathematical
expressions with known curvature from a given library of base functions \cite{RN476}.
DCP is used by the convex optimization modeling languages CVX, CVXPY, Convex.jl, and
CVXR to ensure that the specified optimization problems are convex \cite{RN476}.

In CVXPY, the building blocks such as variables, parameters, numerical constants,
standard arithmetic operators and atomic functions form expressions \cite{RN481}.
DCP analysis will decompose the expression and figure out if the whole expression
is constant, affine, convex or concave. A problem consists of an objective and a list of
constraints. A problem that follows the DCP rules is guaranteed to be solvable by CVXPY.

In fact, the objective functions and constraints functions we use are mostly affine, which is both
convex and concave. It implies that there is extra freedom to change the signs of the
expressions to fit our needs.

There is another rule called disciplined geometric programming (DGP) that we might mention,
but it is much less relevant to our discussions.

\subsection{Disciplined parametrized programming (DPP)}
In CVXPY, Parameters are symbolic representations of constants \cite{RN480}. Disciplined parametrized
programming (DPP) is simply DCP or DGP that has parameters, while satisfying additional restrictions.
Though the problem compilation time is small relative to solution time, DPP can save some time when we
repeatedly try a lot of different parameters, if the formulated problem is only slightly different.

The extra rules for DPP explain why we only define some of the parameters as "Parameter" and the rest
of the relevant numbers are set to constants as the problem is defined.
\begin{itemize}
    \item Under DPP, all parameters are classified as affine, just like variables.
    \item Under DPP, the product of two expressions is affine when at least one of
          the expressions is constant, or when one of the expressions is parameter-affine and the other is parameter-free.
\end{itemize}


\section{SCS}
Splitting Conic Solver (SCS) \cite{scs} is an efficient, open-source solver for large-scale convex quadratic cone problems.
Not only is it widely supported in frameworks such as CVX, CVXPY, YALMIP, Convex.jl and JuMP, it also has useful features like warm starts
and infeasibility detection.
The main reason we select this solver is because it, together with CVXPY, supports complex positive
semidefinite cones in the problem formulation by vectorization \cite{RN482}.
(While other solvers like .. can't or tedious)

\subsection{Supported problem form}
With the assistance of CVXPY, SCS will turn our problem into the following form,
\begin{align}
    \begin{split}\begin{array}{lcr}
            \begin{array}{ll}
                \mbox{minimize}   & (1/2)x^\top  P x + c^\top  x \\
                \mbox{subject to} & Ax + s = b                   \\
                                  & s \in \mathcal{K}
            \end{array}
             & \quad &
            \begin{array}{ll}
                \mbox{maximize}   & -(1/2)x^\top  P x - b^\top  y \\
                \mbox{subject to} & Px + A^\top y + c = 0         \\
                                  & y \in \mathcal{K}^*
            \end{array}
        \end{array}
    \end{split}
\end{align}
where $x \in \mathbf{R}^{n}$ is the primal variable,
$s \in \mathbf{R}^{m}$ is the slack variable,
$A \in \mathbf{R}^{m \times n}$,
$b \in \mathbf{R}^{m}$,
$P \in \mathbf{S}^{n}_{+}$,
$c \in \mathbf{R}^{n}$,
$\mathcal{K} \subseteq \mathbf{R}^{m}$
.

\subsection{Vectorization}
We show below the transformation between a complex positive semidefinite matrix and a vector of real values.
Let $S$ be a $k \times k$ complex matrix
\begin{align}
    S = \begin{bmatrix}
            S_{11} & S_{12} & \cdots & S_{1k} \\
            S_{21} & S_{12} & \cdots & S_{1k} \\
            \vdots & \vdots & \ddots & \vdots \\
            S_{k1} & S_{k2} & \cdots & S_{kk} \\
        \end{bmatrix}
    \in \mathbf{C}^{k \times k}
\end{align}

The vectorization $\text{vec}(S)$ is done such that
$\text{Trace}(Y^\dagger S) = \text{vec}(Y)^\dagger \text{vec}(S)$.
\begin{align}
    \begin{split}\text{vec}(S) = (S_{11}, \sqrt{2} \Re(S_{21}), \sqrt{2} \Im(S_{21}), \ldots, \sqrt{2} \Re(S_{k1}), \sqrt{2} \Im(S_{k1}), S_{22}, \\\sqrt{2}\Re(S_{32}), \sqrt{2}\Im(S_{32}), \dots, S_{k-1,k-1}, \sqrt{2}\Re(S_{k,k-1}), \sqrt{2}\Im(S_{k,k-1}), S_{kk}) \in \mathbf{R}^{k^2}.\end{split}
\end{align}

To recover the complex positive semidefinite matrix from a vectorized solution $s$,
\begin{align}
    \begin{split}\text{mat}(s) =  \begin{bmatrix}
            s_{1}                             & (s_{2} - i s_3) / \sqrt{2}        & \ldots & (s_{2k-2} - is_{2k-1}) / \sqrt{2} \\
            (s_{2} + i s_3) / \sqrt{2}        & s_{2k}                            & \ldots & (s_{4k-5} - is_{4k-4}) / \sqrt{2} \\
            \vdots                            & \vdots                            & \ddots & \vdots                            \\
            (s_{2k-2} + is_{2k-1}) / \sqrt{2} & (s_{4k-5} + is_{4k-4}) / \sqrt{2} & \ldots & s_{k^2}                           \\
        \end{bmatrix}
        ,\end{split}
\end{align}

\subsection{Acceleration for linear system solve}
At each iteration (for example, iteration $k$), SCS solves the following set
of linear equations \cite{linSolve}:
\begin{align}
    \begin{split}\begin{bmatrix}
            R_x + P & A^\top \\
            A       & -R_y   \\
        \end{bmatrix}
        \begin{bmatrix}
            x \\
            y
        \end{bmatrix}
        =
        \begin{bmatrix}
            z^k_x \\
            z^k_y
        \end{bmatrix}\end{split}
\end{align}
for a particular right-hand side $z^k \in \mathbf{R}^{n+m}$
. $R \in \mathbf{R}^{(n+m) \times (n+m)}$ is the diagonal scaling matrix.
The $R_y$ term is negated to make the matrix quasi-definite; we can recover
the solution to the original problem from this modified one. Note that the matrix does not change from
iteration to iteration, which is a major advantage of this approach. Moreover, this system can be solved approximately
so long as the errors satisfy a summability condition, which permits the use of approximate
solvers that can scale to very large problems.


For most of our experiments, we took advantage of the integration of the Intel MKL Pardiso solver in
SCS. The Intel MKL Pardiso solver is a shared-memory multiprocessing parallel direct sparse
solver which is included Intel oneAPI MKL library \cite{linSolve}.

It offers an alternative to the single threaded AMD (approximate minimum degree ordering) / QDLDL libraries which come bundled with SCS.
Pardiso tends to be faster than AMD / QDLDL, especially for larger problems \cite{linSolve}.
Since we are running our experiments
with Intel CPUs, it is very suitable for our needs.

% \subsection{Python interface and arguments}
% Warm start

\section{MOSEK}
MOSEK \cite{mosek} is a high-performance commercial optimization solver designed for convex
optimization problems, including semidefinite programs (SDPs).
It leverages interior-point methods (IPMs), which solve a sequence of barrier
problems to deliver high-precision solutions with robust numerical stability.
This makes MOSEK particularly suitable for applications requiring accurate
primal and dual solutions, such as constraint sensitivity analysis in control theory.

\section{Trade-Off Between SCS and MOSEK}
For solving semidefinite programs (SDPs) in our work, which typically involves
small-scale problems, the trade-off between MOSEK and the Splitting Conic Solver
(SCS) centers on precision, scalability, accessibility, and computational advancements.
MOSEK, utilizing IPMs, achieves high precision with residuals often on the order
of $10^{-8}$ or better, making it ideal for our needs where accurate dual
variables are critical for analyzing constraint importance.
The MOSEK Personal Academic License, freely available to academic users, is
sufficient for our research, providing full access to its capabilities without
additional cost \cite{mosek_academic}.

In contrast, SCS, an open-source solver based on first-order methods like the
Alternating Direction Method of Multipliers (ADMM), offers moderate precision
(residuals of $10^{-4}$ to $10^{-6}$) but requires no license, enhancing its
accessibility.
SCS is designed for scalability, particularly for large-scale problems, and
recent advancements in CVXPY 1.7, released on July 19, 2025, introduced support
for GPU-accelerated solvers, including a GPU-accelerated backend for SCS when
built from source with \texttt{cudss=True} \cite{cvxpy_1_7}.
However, our current work uses CVXPY 1.6, which yet lacks this GPU support.
We will leave the GPU support for future improvements.
For our small-scale SDPs, MOSEK's precision and the sufficiency of its
Personal Academic License outweigh SCS's scalability and license-free benefits,
making MOSEK the preferred choice.
